\section{Introduction}

Language models are increasingly prompted by other language models. Safety evaluations are frequently written by models \citep{perez2022discovering}, automated red-teaming and auditing pipelines generate their own user turns, and in multi-agent systems an orchestrator's instructions to its sub-agents are machine-written. At the same time, models can often tell evaluation transcripts from deployment transcripts \citep{needham2025large}, evaluation awareness is linearly represented and can be steered with behavioural consequences \citep{abdelnabi2025hawthorne,nguyen2025probing,heidari2026evaluation}, and machine-written text is linearly separable from human-written text in LLM activations \citep{chen2025repreguard,quaremba2026linear}. Put together, these facts raise a concrete worry: if a prompt that \emph{reads as LLM-written} shifts a model's safety-relevant behaviour, then results obtained with synthetic prompts may not transfer to deployment with human users, and vice versa.

Testing this is harder than it looks, because ``LLM style'' is bundled with other properties. When an LLM rewrites a human request, the result is typically longer, more explicit, more polite and grammatically cleaner. A behavioural difference between a human prompt and its LLM rewrite could therefore reflect elaboration or politeness rather than perceived authorship, or simply the fact that the text was paraphrased at all. We address this with three kinds of control.

\begin{enumerate}\itemsep0pt
\item \textbf{A factorial rewrite design.} Starting from human-written requests, we create content-matched variants that separate (i) casual-human style from LLM style at matched length, (ii) formality from authorship, via a careful/formal \emph{human-style} variant, (iii) elaboration, via a long LLM ``delegation'' rewrite and a casual-human restyling \emph{of that same long text}, and (iv) paraphrasing as such. All rewrites are checked for content fidelity by a judge model. We mirror the design on natively LLM-written evaluation items.
\item \textbf{Explicit authorship labels.} With the prompt text held fixed, we prepend a note saying that the message was written by a human or generated by an AI, which manipulates \emph{stated} authorship without touching style.
\item \textbf{Activation interventions.} With the prompt text held fixed, we add directions to the residual stream: the ``LLM-vs-casual-human'' style direction, its formality-matched residual, a formality direction, a length direction, a stated-authorship direction, an evaluation-awareness direction, norm-matched random directions and a refusal direction as positive control.
\end{enumerate}

We measure refusal on harmful and borderline-benign requests (JailbreakBench, XSTest), sycophancy towards a user-suggested misconception (TruthfulQA), opinion sycophancy on model-written evaluation items, and answer accuracy, in three open-weight models with activation access (Qwen2.5-7B-Instruct, Llama-3.1-8B-Instruct, Gemma-2-9B-it) and, behaviourally, in three API models (GPT-4.1-mini, Llama-3.3-70B, Claude Haiku 4.5).

\paragraph{Findings.}
\begin{itemize}\itemsep0pt
\item \textbf{The formality-matched authorship effect is null.} Across six models and four behaviours, the contrast between an LLM-style prompt and a careful human-style prompt of the same content and length is small (mostly within $\pm$4 percentage points) and not robustly different from zero (\S\ref{sec:behaviour}). Our design can only detect effects of roughly five points or more per model, and a few isolated contrasts do reach nominal significance.
\item \textbf{What does move behaviour is surface form.} Elaboration (long versus short versions of the same request) lowers refusal by up to 15 points in the open models and GPT-4.1-mini. Casual versus formal phrasing produces smaller, model-specific shifts; for instance Claude Haiku 4.5 declines to answer a model-written political sycophancy item 16--18 points more often in its original or a formal wording than in a casual rewording. LLM style is statistically indistinguishable from formal human style in all of these.
\item \textbf{The internal ``LLM-written'' feature is mostly a formality feature.} A difference-of-means direction separating LLM-style from casual-human prompts reaches AUC $\approx 1.0$ and transfers across datasets, but it does not separate LLM-style from formal-human prompts (AUC $\approx 0.5$) and has cosine $0.67$--$0.89$ with the formality direction. A dedicated probe can still tell LLM-style from formal-human text (AUC $\approx 0.92$), so a residual authorship signal exists, but it is a different and weaker direction.
\item \textbf{Verbalised perception diverges from the internal signal.} GPT-4.1-mini and Claude Haiku 4.5 pick the LLM-style version as ``AI-generated'' about 80\% of the time against a casual version, but are at chance against a formal-human version or the human original. The 7--9B open models are at or below chance; Llama-3.1-8B picks the \emph{casual} version as the AI-generated one 77\% of the time, although its activations separate the two perfectly.
\item \textbf{Steering along the style direction at natural magnitude does little.} Shifting activations by the natural LLM-minus-casual gap, with the text fixed, changes refusal and sycophancy by amounts within the range of norm-matched random directions (\S\ref{sec:steer}), whereas the refusal direction at the same norm moves refusal by tens of points.
\item \textbf{Stated authorship is a different variable.} Labelling a message as AI-generated makes several models somewhat more cautious on borderline requests and less sycophantic, a direction of effect that the style manipulation does not show. The corresponding activation direction is nearly orthogonal to the style direction.
\end{itemize}

For evaluation practice the implication is conditional: synthetic prompts need not change behaviour \emph{because they are synthetic}, but they do change it to the extent that they are more elaborate or more formal than the human prompts they stand in for. Matching those surface properties matters more than disguising authorship.
